%conclusion
%\hl{I have expanded the conclusion to fit IRSs structure. This should be rewritten once the paper is in a more final form}

%\subsection{Summary of work}
In this study, we investigate the drivers and changes in users' posting probability on social networks via a sparse network model and change point detection modelling approach. For the study we collected 3.8 million posts from 502,650 users from WallStreetBets, a subreddit. We examine the impact of individual user behaviour on the Reddit forum Wallstreetbets upon markets. Thereby we provide empirical validation for the theoretical model developed by \citet{Pedersen2022} as well as extending prior research that focuses on aggregated network measures only. 

The DeGroot type model of \citet{Pedersen2022} is a theoretical framework for understanding (meme) stock prices discussed and traded by users of social networks. It being a DeGroot model implies that users of the network change their opinions towards a stock as a result of information they receive from the network, which translates then into price changes. 
Our results demonstrate that changes in users' behaviour significantly explained returns, integrated volatility, and jump volatility, even when controlling for network activity and established metrics measuring influential user impact. We also show that the change in behaviour is not explained by market dynamics, instead it is rooted in the activity on the social network. Including changes in behaviour of users on Reddit into models to explain the market movements, leads to improvements in the modeling fit.
%, leads to an increase in adjusted $R^2$ of $630\%$ for return explainability. %up to $0.45$ for return and $0.8$ for jump explainability, vastly outperforming the competing models. 

Studies often focus upon influential users in networks, but we show that changes in behaviour of less important users explain a larger part of returns, integrated variance and jump volatility than important users do. Influential users change in behaviour results after 6 hours in a significant effect upon market dynamics. This makes sense since a post by a single influential person first needs to reach a credible number of investors before eliciting market reactions. Lesser influential users change in behaviour, if occurring in bulk, ripples instantly and with a various hours lasting effect into the financial markets. Their change in behaviour explains the largest portion of market dynamics, even when controlling for network activity. This finding paves the way towards the construction of variables for stock market movements based on social network users change in behaviour. 

%It has been shown that network activity such as aggregated number of posts at a given time is an important factor in modeling network influence on markets \parencite{Audrino2020, Gan2020, Jiao2020}. While controlling for these and other effects, we study the relationship between the changepoints in the behaviour of individual vital and non-vital users and the market behaviour of GME. We analyse the effect of users change in behaviour on the integrated and jump volatility which we derive from hourly and 5 minute realized volatility, hourly and 5 minute trading volume as well as the corresponding returns. 
%The results show that the change in behaviour of vital and non-vital nodes strongly improves the predictability of said metrics. The adj. $R^2$ from a model without vs. with variables representing behavioural change of individual users, improves by $150\%$. During the market frenzy, these variables contribute towards reaching $0.45$ adj. $R^2$. Likewise for integrated variance, the changepoints of the vital/non-vital nodes excel in explaining market dynamics with the onset of the market frenzy. For jump volatility, even $0.8$ adj. $R^2$ is reached. Even better market predictability is reached when considering to build the models as by trading hour. Given we observe that the developed network model outperforms established metrics for measuring influential user impact, we conclude that the proposed methodology provides a viable framework to investigate a nodes behaviour over time. 

%It is remarkable, that classical network measures such as network activity and network statistics were not able to explain the market movements that well. The explainability of the return dynamics from the changes of users behaviour remained high also after the frenzy, outperforming all other models. This shows that the social network had a longer lasting effect on the markets, going beyond the main market frenzy. 

%\subsection{Theoretical Implications}
%Using our method we address an important gap in the existing literature on the relationship between social networks and stock markets. Our approach focuses on the network micro structure and links individual users' behaviour to stock market behaviour. We show that including this individual change in behaviour results in a larger $R^2$ compared to just aggregated network measures, showing that not only vital nodes in the network but all nodes are important to incorporate. 

Our modelling approaches applicability extends well beyond the scope of the study. With the ever growing number of meme stocks, our method can be used to analyse the impact users on social media have upon market prices. 
While it is extremely challenging to detect change points quickly, our findings suggest that influential users change in behaviour could act as an early warning signal for meme stock market reactions, if identified early enough. We also find that several influential users exhibit a change in behaviour 1-2 months before the market frenzy. %This paves the way towards an early warning signal for meme stock market frenzies. 
Further, our method is applicable to meme coins, which have similar origins as meme stocks. Due to the ease of creating new cryptocurrencies, the market of meme coins is persistently growing and our model provides a method to determine the effect social network users have upon meme coins. 

%\subsection{Practical Contributions}
%Our paper has important managerial implications. First, our findings provide implications for Reddit and other online platforms. Understanding the relationship between online communications on social networks and stock markets helps platforms understand the importance of their role in the economy. Second, meme stocks are having a real impact on funds’ values; however, this influence is extremely unpredictable and volatile. Thus, for fund and portfolio management, a thorough understanding of network evolution and the underlying driving factors is crucial. Our model provides important insights on the relationship between social network discussions and stock prices. Third, a better understanding of the relationship between online social networks and financial markets is crucial for regulators. For these stakeholders, we offer a framework to study the nature of the effect of social networks upon stock markets behaviour.  %within which this phenomenon can be studied. Our research provides guidance to regulators to understand the dynamics on networks and the relationship with the stock market. It shows that we should monitor both the leaders as well as the `crowd' as they can have profound influence upon stock markets. We also 

%\subsection{Limitations and Future Research Opportunities}
This study sets the stage for a number of potential future research opportunities. We show in this study that individual users change in behaviour explains market movements. The GME market frenzy led to the emergence of various meme stocks, with some of them emerging shortly after the market frenzy, such as AMC and BB. A potential future research opportunity would be to study the spillover of the GME frenzy into these assets market frenzies and vice versa. In this study, we distinguish between the relationship of change in behaviour of vital and admissible non-vital users with market movements and show that the studies' results are robust towards different definitions for vital users. An interesting research question is the relationship between individual vital nodes and the admissible non-vital nodes associated with them, which would shed light onto the incremental role single users played in the market frenzy. %network models require variables and structures which suit the network architecture of the social network under investigation. Our study provides a model for the Reddit social network and a different social network would require an extension of the model to its specific architecture. Further, during the study we tested a number of algorithms to detect vital nodes in the network and found them to work not adequately, leading us towards designing a method for vital node identification. Accordingly, future research could design algorithms to identify vital nodes in a dynamic network. 
%As an initial attempt to model and analyze the relationship between individual social network user behavior and stock market analysis this study is prone to several limitations that offer opportunities for future research. First, the current paper focuses on the link structure of a network, an important and interesting extension would be to include the content of posts as well as nodal characteristics to this analysis as \citet{han2020importance} show that the interaction between those two is an important factor determining virality of content in social networks. Second it is challenging to implement our current modeling prospectively due to the manner in which we search change points over a window. Future research should focus on finding changepoint prospectively. Existing appraoches for prospective change detection in networks, such as \cite{zwetsloot2025online}, all focus on aggregated network metrics and ignore the individual nodes. 

%Future research could also explore identification of vital nodes in dynamic networks, we discoverd that existing algorithms did not work adequately as they give us nodes that where very active in only limited amount of time, where influence in a dynamic manner should required an influential node to be active over a longer period of time. Accordingly, future research could design algorithms to identify vital nodes in a dynamic network. 

